\chapter{Selenium}

\section*{What Is This Test?}
The selenium test measures the concentration of selenium in serum, plasma, whole blood, or urine. Selenium is an essential trace element that functions as a component of selenoproteins, most notably the glutathione peroxidases (antioxidant enzymes that reduce hydrogen peroxide and lipid peroxides), thioredoxin reductase, and the iodothyronine deiodinases that activate and inactivate thyroid hormone. Selenium therefore links antioxidant defense, immune function, and thyroid hormone metabolism. Deficiency is uncommon in most populations with adequate soil selenium but occurs in patients on long-term parenteral nutrition without selenium supplementation, in regions with low soil selenium (parts of China, New Zealand, and Scandinavia), and in malabsorption syndromes. Severe deficiency causes Keshan disease (a dilated cardiomyopathy endemic in selenium-deficient regions of China), Kashin-Beck disease (an osteoarthropathy), and impaired immune function. Selenium excess (selenosis) follows excessive supplementation or environmental exposure and causes brittle nails, hair loss, garlic odor, and in severe cases selenosis with neuropathy. Serum selenium is the standard measure of short-term selenium status, while whole-blood and erythrocyte selenium reflect longer-term status.

\section*{Alternative Names}
Serum Selenium; Plasma Selenium; Whole Blood Selenium; Se; Selenoprotein P (related); Glutathione Peroxidase (related)

\section*{Principle}
Selenium is measured by inductively coupled plasma mass spectrometry (ICP-MS), the current reference method, or by hydride generation atomic absorption spectroscopy (HGAAS). In ICP-MS, the sample is digested and ionized in an argon plasma, and selenium isotopes (Se-78, Se-80) are quantified by mass-to-charge ratio. Because argon-based polyatomic interferences (argon-argon dimers) overlap selenium mass peaks, collision-cell or reaction-cell technology (using hydrogen or helium) is used to resolve them. In HGAAS, selenium is reduced to the volatile hydride (H2Se) and measured by atomic absorption, which improves sensitivity by separating selenium from matrix interferences. Specimen collection requires trace-element-free (royal blue-top) tubes to avoid contamination, and hemolyzed samples are rejected. Glutathione peroxidase activity, a functional marker of selenium status, is measured by enzymatic assay. Selenoprotein P, the major selenium transport protein, is measured by immunoassay and reflects selenium supply to tissues.

\section*{History}
Selenium was discovered by Jons Jakob Berzelius in 1817 and named for the Moon (selene). It was considered a toxic element for over a century, causing ``alkali disease'' in livestock. In 1957, Klaus Schwarz and Calvin Foltz demonstrated that selenium prevents liver necrosis in vitamin E-deficient rats, establishing selenium as an essential nutrient. The identification of selenium as a component of glutathione peroxidase by John Rotruck in 1973 explained its antioxidant mechanism. The discovery of Keshan disease (a cardiomyopathy in selenium-deficient regions of China) in the 1970s, and the demonstration that selenium supplementation prevented it, established selenium deficiency as a human disease. Measurement of serum selenium by atomic absorption and later ICP-MS became clinically available in the 1970s--1980s, and selenoprotein research expanded markedly in the 1990s with the identification of more than 25 human selenoproteins.

\section*{Why This Test Is Ordered}
\begin{itemize}
  \item Evaluation of suspected selenium deficiency in patients on prolonged parenteral nutrition without selenium supplementation, which can cause a painful cardiomyopathy (Keshan disease-like) and skeletal muscle weakness
  \item Assessment of patients with malabsorption syndromes: Crohn disease, celiac disease, short bowel syndrome, cystic fibrosis, and chronic pancreatitis
  \item Investigation of unexplained cardiomyopathy in endemic regions or in patients with nutritional risk
  \item Evaluation of patients with thyroid disease or unexplained thyroid function abnormalities, because selenium is required for the iodothyronine deiodinases that convert T4 to T3
  \item Assessment of immune function in malnutrition, HIV, or critical illness, where selenium status predicts outcome
  \item Evaluation of hair loss, nail changes, or skin findings suggesting either deficiency or toxicity
  \item Monitoring of selenium status in patients on high-dose selenium supplements or with occupational or environmental exposure (selenosis)
  \item Assessment in patients with chronic alcohol-related liver disease or end-stage renal disease on dialysis, who are at risk of deficiency
  \item Part of multi-element nutritional assessment (with zinc, copper, magnesium) in malabsorption and refeeding programs
\end{itemize}

\section*{Is This a Primary or Secondary Test}
Serum selenium is a secondary test, ordered when deficiency or toxicity is suspected in at-risk patients. It is not a routine screening test, although it may be included in nutritional assessment panels for malabsorption, parenteral nutrition, and critically ill patients.

\section*{How This Test Is Performed}
A venous blood sample is collected by standard venipuncture into a trace-element-free (royal blue-top) tube. Whole-blood selenium (collected in an EDTA or heparin trace-element tube) is preferred in some circumstances because it reflects longer-term status than serum. Serum is separated and analyzed by ICP-MS or HGAAS. Glutathione peroxidase activity, when ordered, is measured by enzymatic assay on the same specimen. Urine selenium (24-hour collection) is used to assess excretion and exposure. Samples must be protected from contamination and processed promptly. Results are typically available within 3--7 days because trace-element testing is frequently performed at a reference laboratory.

\section*{What Preparation Is Needed}
No fasting is required, though fasting samples are preferred for consistency. Collection must use certified trace-element-free tubes because routine tubes can contaminate the sample. Patients should disclose selenium supplementation and occupational exposure. Because selenium status is affected by inflammation and recent intake, the clinical context (acute illness, recent supplementation) should be documented for interpretation.

\section*{Contraindications and Precautions}
There are no absolute contraindications to standard venipuncture. Important pre-analytical and interpretive cautions: (1) Contamination of the specimen (contact with rubber stoppers, glass, or the skin surface) falsely elevates results; trace-element-free tubes are mandatory. (2) Hemolysis invalidates the sample because red cells concentrate selenium. (3) Serum selenium reflects only recent intake and may be normal despite subclinical tissue deficiency; whole-blood selenium is preferable for longer-term status. (4) The acute-phase response lowers serum selenium during infection, surgery, and critical illness, so low values must be interpreted with inflammatory markers in mind. (5) Selenium deficiency and toxicity have a narrow therapeutic window; supplementation decisions should be guided by both clinical and laboratory assessment. (6) Recent selenium supplements (e.g., 100--200 mcg daily) raise serum levels, so the history of supplementation is essential. (7) Renal failure and hemodialysis lower selenium levels. (8) Because selenium is excreted largely in urine, urine selenium is used specifically to assess excessive exposure.

\section*{Risks}
Minimal risk. Slight pain or bruising at the venipuncture site. Rarely: hematoma, infection, or vasovagal syncope.

\section*{What Happens After the Test}
Results are typically available within 3--7 days. Low selenium in an at-risk patient is treated with oral selenium supplementation (typically 55--200 mcg daily for adults), and the underlying cause (malabsorption, parenteral nutrition, dietary inadequacy) is addressed. For patients on parenteral nutrition, selenium is added to the parenteral solution and levels are monitored periodically. Selenium toxicity is managed by discontinuing supplements and removing the source of exposure; symptoms (nail brittleness, hair loss, garlic breath odor, gastrointestinal distress, and in severe cases neurotoxicity) resolve with cessation. Repeat selenium measurement after 4--12 weeks confirms repletion or normalization.

\section*{Understanding Results}

\subsection*{Below Normal}
\begin{itemize}
  \item \textbf{Selenium deficiency (serum selenium below approximately 70--80 mcg/L, laboratory-dependent):} Causes include prolonged parenteral nutrition without selenium, malabsorption (Crohn disease, short bowel syndrome, cystic fibrosis, celiac disease, chronic pancreatitis), low soil selenium regions, severe malnutrition, and chronic alcohol use.
  \item \textbf{Clinical manifestations:} Keshan disease (dilated cardiomyopathy with heart failure and arrhythmias, historically endemic in selenium-deficient regions of China), skeletal muscle weakness and myopathy, impaired immune function with increased susceptibility to infection, hair and nail changes, and in Kashin-Beck disease (an endemic osteoarthropathy of children), joint and bone changes.
  \item \textbf{In critical illness:} Low selenium is a marker of severity and independently predicts worse outcomes, although randomized trials of selenium supplementation in sepsis have not shown consistent mortality benefit.
  \item \textbf{Interpretive caution:} Serum selenium is falsely lowered by acute-phase inflammation; correlate with CRP and clinical context.
\end{itemize}

\subsection*{Above Normal}
\begin{itemize}
  \item \textbf{Selenium toxicity (selenosis, serum selenium above approximately 150--200 mcg/L, whole blood higher):} Caused by excessive supplementation, high-selenium soils, industrial exposure, or misuse of selenium-containing supplements.
  \item \textbf{Clinical manifestations:} Brittle and shedding nails, hair loss, garlic-like breath odor, metallic taste, gastrointestinal distress, skin rashes, and in severe cases peripheral neuropathy and selenosis with fatigue and irritability.
  \item \textbf{Management:} Discontinue the source; symptoms resolve over weeks to months. Urine selenium measurement documents ongoing exposure.
  \item \textbf{Note:} There is a narrow margin between adequate intake (55 mcg daily) and the tolerable upper limit (400 mcg daily for adults); routine selenium supplementation is not recommended in selenium-replete populations.
\end{itemize}

\section*{Normal Results}
\begin{itemize}
  \item \textbf{Serum selenium (adults):} 70--150 mcg/L (approximately 0.89--1.90 $\mu$mol/L), varying by laboratory, geography, and soil selenium
  \item \textbf{Whole blood selenium:} Approximately 100--200 mcg/L (higher than serum; reflects longer-term status)
  \item \textbf{Plasma selenium:} Similar to serum
  \item \textbf{Urine selenium (24-hour):} 20--90 mcg/day (varies with intake)
  \item \textbf{Glutathione peroxidase activity:} Reference range varies with assay; low activity suggests functional selenium deficiency
  \item \textbf{Daily requirement:} 55 mcg/day for adults (60--70 mcg for women during pregnancy and lactation); tolerable upper intake level 400 mcg/day
  \item \textbf{Conversion:} 1 mcg/L selenium = 0.0127 $\mu$mol/L
\end{itemize}

\section*{Follow-up Tests}
\begin{itemize}
  \item \textbf{Glutathione peroxidase (GPx) activity:} Functional measure of selenium status; low activity confirms deficiency and improves after repletion
  \item \textbf{Other trace elements (zinc, copper, chromium, magnesium):} Deficiencies commonly coexist in malabsorption and parenteral nutrition; measure together for a complete assessment
  \item \textbf{Thyroid function tests (TSH, free T4, free T3):} Selenium is required for deiodinase enzymes; deficiency can impair T4 to T3 conversion
  \item \textbf{Serum albumin and nutritional panel:} Interpret selenium in the context of overall nutritional status
  \item \textbf{Inflammatory markers (CRP or ESR):} A low selenium in the setting of acute-phase inflammation may be spurious
  \item \textbf{Urine selenium (24-hour):} Quantifies excretion and confirms ongoing exposure in suspected toxicity
  \item \textbf{Selenoprotein P:} A transport protein reflecting selenium supply to tissues; useful in research and specialized clinical evaluation
  \item \textbf{Repeat serum selenium after 4--12 weeks of supplementation:} Confirms repletion and guides maintenance dosing
\end{itemize}

\section*{References}
\begin{enumerate}
  \item Rayman MP. Selenium and human health. Lancet. 2012;379(9822):1256--1268.
  \item Weeks BS, Hanna MS, Cooperstein D. Dietary selenium and selenoprotein function. Med Sci Monit. 2012;18(8):RA127--RA132.
  \item Fairweather-Tait SJ, Bao Y, Broadley MR, et al. Selenium in human health and disease. Antioxid Redox Signal. 2011;14(7):1337--1383.
  \item Stoffaneller R, Morse NL. A review of dietary selenium intake and selenium status in Europe and the Middle East. Nutrients. 2015;7(3):1494--1537.
  \item Nuttall KL. Evaluating selenium poisoning. Ann Clin Lab Sci. 2006;36(4):409--420.
\end{enumerate}

\clearpage
